\documentclass[12pt,a4paper]{book}
\usepackage[utf8]{inputenc}
\usepackage[persian]{babel}
\usepackage{amsmath, amssymb, graphicx, pgfplots, booktabs, array, caption}
\usepackage{geometry}
\geometry{left=2.5cm, right=2.5cm, top=2.5cm, bottom=2.5cm}
\pgfplotsset{compat=1.18}

\usepackage{fancyhdr}
\pagestyle{fancy}
\fancyhf{}
\fancyhead[L]{\leftmark}
\fancyfoot[C]{\thepage}
\renewcommand{\headrulewidth}{0.4pt}

\let\oldchapter\chapter
\renewcommand{\chapter}{
    \fancyfoot[C]{}
    \oldchapter
    \thispagestyle{empty}
    \fancyfoot[C]{\thepage}
}

\begin{document}
\pagenumbering{alph}

\begin{titlepage}
    \centering
    \vspace*{2cm}
    {\Huge \textbf{تغذیه درمانی در بیماری‌های مزمن}} \\
    \vspace{0.5cm}
    {\Large \textbf{Nutritional Therapy in Chronic Diseases}} \\
    \vspace{2cm}
    {\large نویسنده: ...............} \\
    \vspace{0.3cm}
    {\large ویراستار: ...............} \\
    \vspace{2cm}
    {\large \today}
\end{titlepage}

\newpage
\thispagestyle{empty}
\vspace*{4cm}
\centering
\textbf{حقوق محفوظ است.} \\
همه‌ی حقوق این کتاب برای نویسنده محفوظ است. \\
هرگونه کپی‌برداری بدون اجازه‌ی کتبی ممنوع است.

\newpage
\thispagestyle{empty}
\vspace*{6cm}
\centering
تقدیم به ...
\vspace*{2cm}

\newpage
\thispagestyle{empty}
\vspace*{6cm}
\centering
سپاس‌گزاری
\vspace*{2cm}

\tableofcontents
\newpage
\listoftables
\newpage
\listoffigures
\newpage

\pagenumbering{arabic}
\setcounter{page}{1}

\chapter{مبانی تغذیه و متابولیسم}
\section{بخش اول: انرژی و نیازهای روزانه}
\section{بخش دوم: درشت‌مغذی‌ها و نقش آنها}
\section{بخش سوم: ریز‌مغذی‌ها و تعادل مایعات}

\begin{table}[h]
\centering
\caption{نیازهای روزانه‌ی انرژی و درشت‌مغذی‌ها}
\label{tab:فصل1}
\begin{tabular}{|c|c|c|c|c|}
\hline
\textbf{سن (سال)} & \textbf{جنسیت} & \textbf{کالری (کیلوکالری)} & \textbf{پروتئین (گرم)} & \textbf{کربوهیدرات (گرم)} \\
\hline
۱۸-۳۰ & مرد & ۲۴۰۰ & ۵۶ & ۳۳۰ \\
۱۸-۳۰ & زن & ۲۰۰۰ & ۴۶ & ۲۶۰ \\
۳۱-۵۰ & مرد & ۲۲۰۰ & ۵۴ & ۳۰۰ \\
۳۱-۵۰ & زن & ۱۸۰۰ & ۴۴ & ۲۴۰ \\
۵۱+ & مرد & ۲۰۰۰ & ۵۰ & ۲۷۰ \\
۵۱+ & زن & ۱۶۰۰ & ۴۰ & ۲۲۰ \\
\hline
\end{tabular}
\end{table}

\begin{figure}[h]
\centering
\begin{tikzpicture}
\begin{axis}[
    xlabel={گروه غذایی},
    ylabel={درصد مصرف},
    title={هرم غذایی: توزیع درصدی گروه‌های غذایی},
    symbolic x coords={کربوهیدرات,سبزیجات,میوه‌ها,پروتئین,چربی‌ها},
    xtick=data,
    ybar,
    ymin=0,
    ymax=40,
    nodes near coords,
    nodes near coords align={vertical},
]
\addplot coordinates {(کربوهیدرات,30) (سبزیجات,25) (میوه‌ها,20) (پروتئین,15) (چربی‌ها,10)};
\end{axis}
\end{tikzpicture}
\caption{هرم غذایی}
\label{fig:فصل1}
\end{figure}

\chapter{تغذیه در دیابت و چاقی}
\section{بخش اول: کنترل قند خون و شاخص گلیسمی}
\section{بخش دوم: رژیم‌های کاهش وزن و مدیریت چاقی}

\begin{table}[h]
\centering
\caption{شاخص گلیسمی (GI) و بار گلیسمی (GL) مواد غذایی}
\label{tab:فصل2}
\begin{tabular}{|c|c|c|c|}
\hline
\textbf{ماده غذایی} & \textbf{GI} & \textbf{GL} & \textbf{نوع} \\
\hline
نان سفید & ۷۰ & ۱۰ & بالا \\
سیب & ۳۶ & ۶ & پایین \\
برنج قهوه‌ای & ۵۰ & ۱۶ & متوسط \\
سیب‌زمینی & ۸۵ & ۲۶ & بالا \\
عدس & ۲۹ & ۵ & پایین \\
\hline
\end{tabular}
\end{table}

\begin{figure}[h]
\centering
\begin{tikzpicture}
\begin{axis}[
    xlabel={زمان (دقیقه)},
    ylabel={قند خون (میلی‌گرم/دسی‌لیتر)},
    title={منحنی قند خون پس از مصرف مواد غذایی},
    xmin=0,
    xmax=120,
    ymin=80,
    ymax=180,
    legend pos=north west,
    grid=both,
]
\addplot[thick, red] coordinates {(0,90) (30,160) (60,170) (90,120) (120,95)};
\addlegendentry{GI بالا}
\addplot[thick, blue] coordinates {(0,90) (30,110) (60,130) (90,120) (120,100)};
\addlegendentry{GI پایین}
\end{axis}
\end{tikzpicture}
\caption{منحنی قند خون}
\label{fig:فصل2}
\end{figure}

\chapter{تغذیه در بیماری‌های قلبی-عروقی}
\section{بخش اول: چربی‌های خون و عوامل خطر}
\section{بخش دوم: رژیم مدیترانه‌ای و سلامت قلب}
\section{بخش سوم: نقش سدیم و فشار خون}
\section{بخش چهارم: آنتی‌اکسیدان‌ها و سلامت عروق}

\begin{table}[h]
\centering
\caption{ترکیب اسیدهای چرب در روغن‌ها}
\label{tab:فصل3}
\begin{tabular}{|c|c|c|c|c|}
\hline
\textbf{روغن} & \textbf{اشباع (\%)} & \textbf{تک‌غیراشباع (\%)} & \textbf{چند‌غیراشباع (\%)} & \textbf{ترانس (\%)} \\
\hline
زیتون & ۱۴ & ۷۳ & ۱۱ & ۰ \\
نارگیل & ۸۶ & ۶ & ۲ & ۰ \\
آفتاب‌گردان & ۱۰ & ۲۰ & ۶۶ & ۰ \\
کره & ۶۳ & ۲۶ & ۴ & ۳ \\
مارگارین & ۳۰ & ۴۰ & ۲۰ & ۱۰ \\
\hline
\end{tabular}
\end{table}

\begin{figure}[h]
\centering
\begin{tikzpicture}
\begin{axis}[
    xlabel={نوع رژیم غذایی},
    ylabel={میزان (میلی‌گرم/دسی‌لیتر)},
    title={مقایسه کلسترول و تری‌گلیسرید در دو رژیم},
    symbolic x coords={غربی,مدیترانه‌ای},
    xtick=data,
    ybar=0.5cm,
    ymin=0,
    ymax=250,
    legend pos=north west,
    nodes near coords,
    nodes near coords align={vertical},
]
\addplot coordinates {(غربی,220) (مدیترانه‌ای,170)};
\addplot coordinates {(غربی,180) (مدیترانه‌ای,130)};
\legend{کلسترول,تری‌گلیسرید}
\end{axis}
\end{tikzpicture}
\caption{مقایسه رژیم‌های غذایی}
\label{fig:فصل3}
\end{figure}

\chapter{تغذیه در بیماری‌های کلیوی و سرطان}
\section{بخش اول: محدودیت‌های غذایی در بیماری کلیوی}

\begin{table}[h]
\centering
\caption{محدودیت‌های غذایی در مراحل CKD}
\label{tab:فصل4}
\begin{tabular}{|c|c|c|c|}
\hline
\textbf{مرحله CKD} & \textbf{پروتئین (گرم/کیلوگرم)} & \textbf{پتاسیم (میلی‌گرم)} & \textbf{فسفر (میلی‌گرم)} \\
\hline
مرحله ۱-۲ & ۰.۸ & ۲۰۰۰ & ۸۰۰ \\
مرحله ۳ & ۰.۶ & ۱۵۰۰ & ۶۰۰ \\
مرحله ۴ & ۰.۴ & ۱۰۰۰ & ۵۰۰ \\
مرحله ۵ (دیالیز) & ۱.۲ & ۲۰۰۰ & ۱۰۰۰ \\
\hline
\end{tabular}
\end{table}

\begin{figure}[h]
\centering
\begin{tikzpicture}
\begin{axis}[
    xlabel={مراحل CKD},
    ylabel={درصد بقا},
    title={نمودار بقای بیماران کلیوی بر اساس وضعیت تغذیه},
    symbolic x coords={مرحله۱,مرحله۲,مرحله۳,مرحله۴,مرحله۵},
    xtick=data,
    ybar,
    ymin=0,
    ymax=100,
    nodes near coords,
    nodes near coords align={vertical},
]
\addplot coordinates {(مرحله۱,90) (مرحله۲,75) (مرحله۳,55) (مرحله۴,30) (مرحله۵,10)};
\end{axis}
\end{tikzpicture}
\caption{بقای بیماران کلیوی}
\label{fig:فصل4}
\end{figure}

\end{document}
